\section{Introduction}
\label{sec:intro}

HIV-1 broadly neutralizing antibodies (bnAbs) are specialized immune system
proteins capable of neutralizing a wide range of viral variants by binding to
comparatively conserved regions on the viral envelope
glycoprotein~\cite{sok2018recent,liu2020bnabs}. Unlike ordinary antibodies,
bnAbs develop over years through a process called somatic hypermutation, where
B-cell receptor genes undergo iterative mutations and selection for increasing
binding affinity.[add cite] This process yields complex clonal lineages
(Figure~\ref{fig:pgt121_lineage} shows an example maturation pathway)—groups of
antibodies that share a common genetic ancestor (the Unmutated Common Ancestor,
or UCA) but diverge in sequence and function as they mature. Predicting
neutralization breadth (the proportion of viral strains neutralized) directly
from the antibody's variable heavy-chain ($V_H$) sequence is challenging because
the functional target depends on a highly diverse sequence, and laboratories
measure it using inconsistent virus panels[add cite].

\begin{figure}[t]
  \centering
  \includegraphics[width=\linewidth]{figures/pgt121_lineage.png}
  \caption{Clonal lineage maturation tree for the PGT121 family of broadly neutralizing antibodies, illustrating somatic hypermutation paths starting from the germline precursor (UCA) through ancestral intermediates to mature variants (reproduced from Sok~\etal~\cite{Sok2013}).}
  \label{fig:pgt121_lineage}
\end{figure}


Protein language models (PLMs) utilize transformer architectures to learn
contextual, high-dimensional sequence representations from millions of unlabeled
proteins. The state-of-the-art multimodal PLM ESM3 learns through self-supervision
to reason over protein sequence, 3D structure, and
function~\cite{hayes2025esm3,evolutionaryscale2024esm3open}. This raises two
questions. First, can the general-protein representations that ESM3 learns
support out-of-lineage bnAb modeling (predicting properties for unseen antibody
families) using only the heavy-chain sequence? Second, does unsupervised domain
adaptation to antibody repertoires (using sequence-only data from the Observed
Antibody Space [OAS] database~\cite{olsen2022oas}) improve this representation
transfer?

We guard against data leakage because antibody datasets
contain many somatically related variants; thus, a naive
random split can place members of the same clonal lineage in both training and
validation folds. This allows models to train on ancestors or descendants,
rather than learning generalizable properties. In
this project, we implement a metadata-derived lineage grouping for
lineage isolation across folds. Because the
study uses the same dataset to compare multiple model architectures (heads) and pooling
methods, the cross-validation score represents model-selection
estimates.

We organize the study around three contributions. First, we explain the
metadata-derived lineage grouping and group-resampled uncertainty
analysis for classification[is this novel???] and regression on the Compile/Analyze Tool for
Neutralization Assay Data (CATNAP) database~\cite{yoon2015catnap}. Second, we
present a matched evaluation of sequence features, frozen ESM3, and OAS-tuned
ESM3 (via LoRA) under identical downstream folds and heads. Third, we perform
a controlled analysis of VAE interpolation on residue-resolution sequence embeddings.
This separates observations from unsupported manifold claims.
